\section{绪论}

\subsection{研究背景}

视觉语言模型在近年取得了长足进展。CLIP 等对比式模型通过大规模预训练将图像和文本映射到统一语义空间，BLIP-2、LLaVA 和 Qwen-VL 等生成式模型进一步将视觉编码器与大语言模型连接，实现了图像描述、视觉问答等能力。这些模型的共同局限在于：当输入图像出现低分辨率、模糊、噪声或遮挡等退化时，语义理解和文本生成质量会显著下降。如何在退化条件下保持视觉语言模型的鲁棒性，是将其推向实际应用需要面对的关键挑战。

另一方面，EEG 技术能以毫秒级分辨率记录视觉刺激引发的脑电活动。近年来，从 EEG 中解码视觉语义信息的研究取得了积极进展——已有多项工作证明 EEG 信号中包含可解码的物体类别和语义嵌入信息。这自然引发我们的关注：当视觉传感器（摄像头）因图像退化而输出不可靠的特征时，人脑这个天然视觉系统是否能为机器提供补充信息？正是带着这个想法，我们开展了本次课程设计。

需要说明的是，我们的研究定位不是从脑电中直接解码语言，而是将 EEG 作为一种视觉语义增强信号，在所有实验结论中通过与 shuffled EEG 和 random EEG 的对照来保证判断的客观性。

\subsection{课程设计目标与意义}

本次课程设计的目标是在 Thought2Text 小样本 EEG-image 数据集上，构建视觉退化条件下的 EEG 辅助视觉语言理解系统，并通过系统的对照实验评估真实 EEG 的语义辅助效果。

从实践训练的角度，该课题让我们综合运用了深度学习模型设计、预训练模型适配、多模态融合、大模型参数高效微调（LoRA）和对比学习（InfoNCE）等多项技术，在一个真实数据集上完成了从问题调研到最终评估的完整流程。从方法验证的角度，通过 real EEG、shuffled EEG、random EEG 和 vision-only 的四重对照设计，我们试图为``EEG 是否携带可机器利用的视觉语义信息''这一开放问题提供可验证的实证参考。

\subsection{研究内容与技术路线}

本次课程设计完成了以下核心工作：

\begin{enumerate}
  \item 搭建了视觉退化条件下的 EEG 辅助语义识别与生成式描述实验框架，定义六种退化条件和四种 EEG 对照模式，采用 image-level split 防止数据泄漏；
  \item 设计并实现 A2 时间-频谱-空间 EEG 语义融合模型，从 EEG 的时间、频谱和空间通道三个维度并行编码，通过门控残差融合对退化图像进行语义补充；
  \item 实现 VTF token 级融合方法，将 EEG token 通过交叉注意力注入 CLIP ViT visual token，生成增强视觉 token，为后续接入生成式模型做准备；
  \item 构建基于 Qwen2-VL 的生成式图像描述模块，使用 LoRA 高效微调和多候选重排序，对比五种技术方案的优劣，产出自由形式 caption。
\end{enumerate}

在技术路线上，我们采用了从简到繁、先验证后扩展的策略：先用 CLIP 预训练语义空间作为跨模态桥梁降低数据量要求，再在分类任务上验证 EEG 语义增强的有效性，确认后再逐步推进到 token 级融合和生成式输出。这种分层推进的策略让每一步都有可检查的中间结果。

\subsection{报告结构}

本报告共分十章。第 1 章绪论后，第 2 章介绍相关技术基础；第 3--4 章介绍数据集、任务和系统总体设计；第 5--7 章分别详述 A2、VTF 和生成式 EVLM 三个核心模块的方法与实现；第 8 章给出完整实验结果与分析；第 9 章概述工程实现；第 10 章总结与展望。
